\documentclass[11pt,letterpaper]{article}

\usepackage{graphicx}           % Graphics package
\usepackage{amsmath} 
\usepackage{mathrsfs}
\usepackage[colorlinks=true]{hyperref} 
\usepackage{color}
\usepackage{amsfonts}
\usepackage[textheight=9in, textwidth=7.5in, letterpaper]{geometry}
\usepackage[makeroom]{cancel}
\usepackage{cite}
\usepackage{sectsty}
\usepackage{empheq}
\usepackage{appendix}
\usepackage{enumerate} 
\usepackage{simpler-wick}
\usepackage{amssymb}

\sectionfont{\fontsize{12}{12}\selectfont}
\subsectionfont{\fontsize{12}{12}\selectfont}

\newcommand*\widefbox[1]{\fbox{\hspace{2em}#1\hspace{2em}}}
\newcommand{\LP}{\left(}
\newcommand{\RP}{\right)}
\newcommand{\mc}[1]{\mathcal{#1}}
\newcommand{\R}{\mathbb{R}}
\newcommand{\C}{\mathbb{C}}
\newcommand{\Z}{\mathbb{Z}}
\newcommand{\D}{\partial}
\newcommand{\KET}[1]{\left| #1 \right\rangle }
\newcommand{\BRA}[1]{\left\langle #1 \right| }
\newcommand{\IP}[2]{\left\langle #1 \middle| #2 \right\rangle}
\newcommand{\PA}[2]{\frac{\partial #1}{\partial #2}}



%%%%%%%%%%%%%%%%---------------------MOST USEFUL COMMAND EVER---------------------%%%%%%%%%%%%%%%%%%%%%%
\newcommand{\fixme}[1]{{\bf {\color{red}[#1]}}}
\newcommand{\draftmode}{\usepackage[notref,notcite]{showkeys}}
\providecommand*\showkeyslabelformat[1]{\normalfont\sffamily\footnotesize#1}

%%%%%%%%%%%%%%%


\setcounter{tocdepth}{4}
\setcounter{secnumdepth}{4}

\begin{document}

\title{Multipole Moments}

\author{Mathew Calkins}

\date{\today}

 
\maketitle

\abstract{Extremely rough notes on multipole moments. Maybe these will turn into something coherent, maybe not!}

\tableofcontents




\section{Spherical Harmonics}
Spherical harmonics most naturally arise when solving the 3D Euclidean Laplace equation \begin{equation}
\nabla \phi = 0, \quad \phi: \R^3 \to \R
\end{equation}
by the ansatz of separation of variables in spherical coordinates \begin{equation}
\phi(r, \theta, \phi) = R(r) Y(\theta, \phi)
\end{equation}
whereupon the angular function equivalently obeys \begin{equation}
- \lambda = \frac{1}{Y} \frac{1}{\sin \theta} \D_\theta \LP \sin \theta \D_\theta Y \RP + \frac{1}{Y} \frac{1}{\sin^2 \theta} \D_\phi^2 Y
\end{equation}
for $\lambda$ a constant also appearing in the ODE for $R$. Further separating \begin{equation}
Y(\theta, \phi) = \Theta(\theta) \Phi(\phi)
\end{equation}
gives the system \begin{equation}
\begin{aligned}
\Phi'' & = - m^2 \Phi, \\
\lambda \sin^2 \theta + \frac{\sin \theta}{\Theta} \frac{d}{d \theta} \LP \sin \theta \frac{d \Theta}{d \theta} \RP & = m^2
\end{aligned}
\end{equation}
for $m \in \R$ an additional constant. The first equation forces $m \in \Z$. The second, plus demanding regularity at the poles $\theta = 0, \pi$, forces \begin{equation}
\lambda = \ell(\ell + 1) \mbox{ for } \ell \in \Z, \ell \geq |m|.
\end{equation}
So we have a basis of solutions \begin{equation}
Y_\ell^m (\theta, \phi) \propto e^{i m \phi} P_\ell^m (\cos \theta), \quad m = -\ell, - \ell + 1, \ldots, \ell
\end{equation}
where $P_\ell^m$ are the associated Legendre polynomials. In that case, for summations over basis solutions converging at large $r$, the radial function is then \begin{equation}
R(r) \propto \frac{1}{r^{\ell + 1}}.
\end{equation}
We can say more about the functions $Y$, namely that they make up exactly the complete set of  eigenfunctions of the spherical Laplacian \begin{equation}
r^2 \nabla^2 Y_\ell^m = - \ell(\ell + 1) Y_\ell^m.
\end{equation}
This notation is given in a slightly ambiguous way in some presentation, leaving it unclear where we mean the 3D Euclidean Laplacian applied to a function which just happens to be $r$-independent, or the Laplacian intrinsic to $S^2$ applied to a map $S^2 \to \R$? As it turns out, these agree even in higher dimensions. To see this, take the usual curvilinear expression for the Laplacian of a scalar \begin{equation}
\nabla^2 f = \frac{1}{\sqrt{g}} \D_a \LP \sqrt{g} g^{ab} \D_b f \RP
\end{equation}
and apply this to the Euclidean metric for $\R^N$ written in generalized spherical coordinates \begin{equation}
ds^2 = dr^2 + r^2 d \Omega^2 \implies \sqrt{g} = r^{N-1}
\end{equation}
Considering the cases where $a,b$ are both angles, both radial, or a mixture, we find that \begin{equation}
\nabla_{\R^N}^2 f = \frac{1}{r^{N-1}} \PA{}{r} \LP r^{N-1} \PA{f}{r} \RP + \frac{1}{r^2} \Delta^2_{S^{N-1}} f.
\end{equation}
So a function $f: S^2 \to \R$ is annihilated by the sphere Laplacian iff, when interpreted as a radius-independent function on $\R^N$, it is annihilated by the Euclidean Laplacian.\\
\\
\
We have just diagonalized the linear operator\footnote{As usual in theoretical physics, here we are assuming as much smoothness from our test functions as we could ever plausibly need. The domains and codomains here are kind of cartoon placeholders, and a real treatment of representation theory will give the right ones.} \begin{equation}
\nabla^2_{\R^3}: C^\infty(S^2, \C) \to C^\infty(S^2, \C).
\end{equation}
If we interpret its domain and codomain as linear spaces endowed with the unique-up-to-scaling rotation-invariant inner product \begin{equation}
\langle f, g \rangle \propto \int_{S^2} d \Omega \bar{f(\Omega)} g(\Omega)
\end{equation} 
then $\nabla^2_{\R^3}$ becomes a Hermitian operator, so its eigenfunctions form a complete orthonormal basis. In fact, since we can take infinite sums over this space of differentiable functions and get non-differentiable limits, they span the larger space \begin{equation}
L^2(S^2, \C) = \mbox{Span}_\C \LP Y_\ell^m \RP_{|m| \leq \ell}.
\end{equation}
Again, refer to a real math text to get the careful version.



\section{$SO(3)$ Irreps}
Recall that $SO(3)$ has a distinct\footnote{Up to the usual conjugation equivalence.} complex irrep $V_\ell$ of every odd dimension\footnote{If we loosen to representations of $so(3)$ that don't exponentiate to single-valued reps of the group but generally represent the double cover $SU(2)$, we get further irreps at half-integer $\ell \in \frac{1}{2} + \Z_{\geq 0}$.} $2 \ell + 1$ for $\ell \in \Z_{\geq 0}$. These are generated by the Lie algebra basis elements which, in one standard choice, have matrix elements \begin{equation}
\begin{aligned}
[L^{(\ell)}_x]_{ij} & = \frac{1}{2} (\delta_{i, j+1} + \delta_{i+1,j}), \sqrt{(\ell + 1) (i + j - 1) - ij} \\
[L^{(\ell)}_y]_{ij} & = \frac{1}{2i} (\delta_{i, j + 1} - \delta_{i + 1, j}) \sqrt{(\ell + 1)( i + j - 1) -ij }, \\
[L^{(\ell)}_z]_{ij} & = (\ell + 1 - j) \delta_{ij}.
\end{aligned}
\end{equation}
Each irrep is labeled by the eigenvalues of the unique casimir \begin{equation}
\vec{L}^2 \equiv L_x^2 + L_y^2 + L_z^2 = - \ell(\ell + 1) \mathbf{1}_{2 \ell + 1}
\end{equation}
with degeneracy within each irrep broken by (in the usual choice) $L_z$, having eigenvalues $m \in \Z$, $|m| \leq \ell$.\\
\\
\
$SO(3)$ has a natural representation on the space of square-integrable maps $S^2 \to \C$, acting by pullback as \begin{equation}
f \xmapsto{R \in SO(3)} f \circ R.
\end{equation}
Under this choice of representation, $L^2(S^2, \C)$ should decompose as the direct sum of irreps \begin{equation}
L^2(S^2, \C) \simeq \bigoplus_{\ell = 0}^\infty n_\ell V_\ell
\end{equation}
for multiplicies $n_\ell \in \Z_{\geq 0}$. On the other hand, we see geometrically that each eigenspace \begin{equation}
\mbox{Span}_\C \LP Y_\ell^m \RP_{|m| \leq \ell}^{\mbox{\tiny{fixed }} \ell} = \{ f: S^2 \to \C \mbox{ with } \nabla^2_{S^2} f = - \ell (\ell + 1) \} \subset L^2(S^2, \C)
\end{equation}
is closed under rotation, so we already have a decomposition into closed subspaces \begin{equation}
L^2(S^2, \C) = \bigoplus_{\ell = 0}^\infty \mbox{Span}_\C \LP Y_\ell^m \RP_{|m| \leq \ell}^{\mbox{\tiny{fixed }} \ell}.
\end{equation}
The irrep decomposition must be either this decomposition or some strict refinement of the same. Since the dimensions here at each $\ell$ are $2 \ell + 1$, which matches the dimension of $V_\ell$, it is tempting to guess that \begin{equation}
\mbox{Span}_\C \LP Y_\ell^m \RP_{|m| \leq \ell}^{\mbox{\tiny{fixed }} \ell} \stackrel{?}{\simeq} V_\ell.
\end{equation}
This turns out to be true! We can see this by an exact calculation. In this pullback representation, in the physics convention \begin{equation}
\mbox{(group element)} = e^{i \mbox{\tiny{(algebra element)}}}
\end{equation}
the standard basis for the $so(3)$ is represented by the differential operators \begin{equation}
L_i = - i \epsilon_{ijk} x_j \D_k.
\end{equation}
A concrete calculation gives the unique Casimir here as \begin{equation}
\vec{L}^2 = L_x^2 + L_y^2 + L_z^2 = \nabla^2_{\R^3}
\end{equation}
so our unique Casimir acts on level-$\ell$ harmonics as \begin{equation}
\vec{L}^2 Y_\ell^m = - \ell(\ell + 1) Y_\ell^m
\end{equation}
and up to isomorphism our level-$\ell$ harmonics all live in $V_\ell$. Next, we can compute exactly \begin{equation}
L_z = - i (x \D_y - y \D_x) = \D_\phi \implies L_z Y_\ell ^m = m Y_\ell^m.
\end{equation}
so every value of $m$ in $V_\ell$ is accounted for. So we have exactly \begin{equation}
\mbox{Span}_\C \LP Y_\ell^m \RP_{|m| \leq \ell}^{\mbox{\tiny{fixed }} \ell} \simeq V_\ell.
\end{equation} 







\section{Harmonic Polynomials}
Denote by $\mbox{Hom}_\ell(\R^3)$ the space of all degree-$\ell$ homogeneous polynomial functions from $\R^3$ to $\C$, and by $\mc{H}_\ell(\R^3)$ the space of all harmonic polynomial functions from $\R^3$ to $\C$ which are homogeneous of degree $\ell$. \begin{equation}
\mc{H}_\ell(\R^3) = \{ p \in\mbox{Hom}_\ell(\R^3) \mbox{ with } \Delta p = 0 \}.
\end{equation}
For example, at $\ell = 0$ we get the space of all linear functions $\R^3 \to \C$. \begin{equation}
\mc{H}_1(\R^3) = \mbox{span}_\C (x_1, x_2, x_3).
\end{equation}
At $\ell = 2$ we get the slightly more involved space \begin{equation}
\begin{aligned}
\mc{H}_2(\R^3) &  = \{ a_{11} x_1^2 + a_{22} x_2^2 + a_{33} x_3^2 a_{12} x_1 x_2 + a_{13} x_1 x_3 + a_{23} x_2 x_3 \mbox{ such that } x_1 + x_2 + x_3 = 0\} \\
& = \mbox{span}_\C \LP x_1 x_2, x_1 x_3, x_2 x_3, x_1^2 - x_2^2, x_1^2 - x_3^2 \RP. 
\end{aligned}
\end{equation}
How large is $\mc{H}_\ell(\R^3)$ in general? The space of degree-$\ell$ polynomials finds a natural basis in the space of monomials \begin{equation}
\mbox{Hom}_\ell(\R^3) = \mbox{Span}_\C \LP \{ x_1^{a_1} x_2^{a_2} x_3^{a_3} \mbox{ with } a_1 + a_2 + a_3 = \ell \} \RP.
\end{equation}
Specifying such a monomial means specifying $a_1 \in \{0, 1, \ldots, \ell\}$, then specifying $a_2 \in \{\ell - a_1, \ell - a_1 + 1, \ldots, \ell\}$, after which $a_3$ is fixed. So the number of such choices is exactly \begin{equation}
\begin{aligned}
\dim_\C \LP \mbox{Hom}_\ell(\R^3) \RP & = \sum_{i=0}^\ell \sum_{j=\ell - i}^{\ell} 1 \\
& = \sum_{i=0}^\ell (i+1) \\
& = \frac{(\ell + 2)(\ell + 1)}{2}.
\end{aligned}
\end{equation}
We can compute the dimension of $\mc{H}_\ell(\R^3)$ by observing that the restriction of the Laplacian to arguments of degree $\ell$ \begin{equation}
\Delta_\ell : \mbox{Hom}_\ell(\R^3) \to \mbox{Hom}_{\ell - 2}(\R^3), \quad \ell \geq 2
\end{equation}
is surjective. As a low-order non-trivial example, $x^2 y^2$ can be written non-uniquely as the Laplacian \begin{equation}
\begin{aligned}
xy & = \Delta \LP a x^4 + b x^3 y + c x^2 y^2 + d x y^3 + e y^4 \RP \\
& = 12 a x^3 + 6 b x y + 2 c (x^2 + y^2) + 6 d xy + 12 y^2
\end{aligned}
\end{equation}
which has solutions $6(b+d) = 1$ with all other coefficients vanishing. So rank-nullity gives the dimension as \begin{equation}
\begin{aligned}
\dim_\C \LP \mc{H}_\ell(\R^3) \RP & = \dim_\C \LP \mbox{ker} \ \Delta_\ell \RP \\
& = \dim_\C \LP \mbox{Hom}_\ell(\R^3) \RP - \dim_\C \LP \mbox{Hom}_{\ell - 2} (\R^3) \RP \\
& = \frac{(\ell + 2)(\ell + 1)}{2} - \frac{\ell(\ell - 1)}{2} \\
& = 2 \ell + 1.
\end{aligned}
\end{equation}
We introduced these homogeneous harmonic polynomials just now because there is a cute trick for turning them into a spherical harmonic. First, observe that a degree-$\ell$ homogeneous polynomial $p: \R^3 \to \C$ gives rise in a canonical way to $P: S^2 \to \R$ as \begin{equation}
p(\vec{x}) = r^\ell p(\hat{x}) \equiv r^\ell p(\Omega) \equiv r^\ell P(\Omega) , \quad \vec{x} = r \hat{x}.
\end{equation}
In one of several equivalent phrasings, $P$ is the restriction of $p$ to $S^2$. Now if $p$ is furthermore harmonic w.r.t. the Euclidean Laplacian, we can compute explicitly \begin{equation}
\begin{aligned}
0 & = \nabla^2_{\R^3} p (\vec{x}) \\
& = \frac{1}{r^2} \PA{}{r} \LP r^2 \PA{(r^\ell P(\Omega))}{r} \RP + \frac{1}{r^2} \Delta^2_{S^2} (r^\ell P(\Omega)) \\
& = \frac{P(\Omega)}{r^2} \PA{}{r} \LP r^2 \PA{(r^\ell )}{r} \RP + r^{\ell - 2} \Delta^2_{S^2} P(\Omega) \\
& = \ell (\ell + 1) r^{\ell - 2} P(\Omega) + r^{\ell - 2} \Delta^2_{S^2} P(\Omega) \\
\implies \Delta^2_{S^2} P(\Omega) & = - \ell (\ell + 1) P(\Omega).
\end{aligned}
\end{equation}
So $P$ is an eigenfunction of the sphere Laplacian, hence is a linear combination of the level-$\ell$ spherical harmonics. So we have cooked up a linear map \begin{equation}
\mc{H}_\ell(\R^3) \to \mbox{Span}_\C \LP Y_\ell^m \RP_{|m| \leq \ell}, \quad p \mapsto p|_{S^2} \equiv P.
\end{equation}
Some elementary algebraic geometry tells us this map is injective: if two polynomials $p,q: \R^3 \to \C$ of the same degree $\ell$ give rise to identical restrictions $P, Q: S^2 \to \C$, then their difference $p - q: \R^3 \to \C$ vanishes on the locus $x^2 + y^2 + z^2 - 1 = 0$. \begin{equation}
(p - q)(\vec{x}) = 0 \mbox{ whenever } |\vec{x}| = 1.
\end{equation}
But $p,q$ have the same degree-$\ell$ scaling behavior by assumption, hence so does their difference, and for general $\vec{x} = r \hat{x} \in \R^3$ we have \begin{equation}
(p - q)(\vec{x}) = r^\ell (p - q).
\end{equation}
So $p = q$, and as claimed our restriction map is injective. On the other hand, we have already computed the dimension of the domain and codomain of this map, finding $2 \ell + 1$ in each case. So our injective map must be an isomorphism, and we conclude that we have furnished a linear equivalence between the space of degree-$\ell$ harmonic polynomials and the degree-$\ell$ spherical harmonics when both are taken as complex vector spaces. \begin{equation}
\therefore \mc{H}_\ell(\R^3) \simeq \mbox{Span}_\C \LP Y_\ell^m \RP_{|m| \leq \ell}
\end{equation}



\section{STF tensors}
We find a further equivalence with the space of symmetric trace-free (STF) tensors over $\C^3$. To be clear, here we mean literally tensors over the linear space $\C^3$, not sections of tensor bundles fibered over $\C^3$ taken as a manifold. We also mean specifically contravariant tensors, which in the standard basis have exclusively lowered indices. To furnish this equivalence, observe that to any rank-$\ell$ contravariant tensor $T$ over $\C^3$ we can associate in a canonical way the degree-$\ell$ homogeneous polynomial \begin{equation}
P_T \equiv P_{a_1 \cdots a_\ell} x_{a_1} \cdots x_{a_\ell}.
\end{equation}
This contraction is sensitive only to the part of $P$ which is totally symmetric in all indices, so restricting our attention to such symmetric tensors is exactly the restriction that makes this map injective without cutting down its image. Better yet, since every polynomial can be written in this way, this mapping is a bijection. This polynomial has Euclidean Laplacian \begin{equation}
\begin{aligned}
\nabla^2_{\R^3} P_T & = g^{ab} \D_{x_a} \D_{x_b} P_{a_1 \cdots a_\ell} x_{a_1} \cdots x_{a_\ell} \\
& = \ell(\ell + 1) g^{ab} P_{a b a_3 \cdots a_\ell}
\end{aligned}
\end{equation}
so $P_T$ is harmonic iff $P$ is trace-free\footnote{Note that since $P$ is totally symmetric, there is no ambiguity in which pairing of indices we are using for the trace; any one pairing vanishes iff every pairing vanishes.}. So we have yet another equivalence \begin{equation}
P_T \in \mc{H}_\ell(\R^3) \leftrightarrow P \in \mbox{Sym}_0^\ell(\C^3) \equiv \{ \mbox{symmetric traceless rank-$\ell$ polynomials over $\C^3$}\}.
\end{equation}


\section{Multipole Moments}
Legendre polynomials arise from from large-distance moments of inverse-distance potentials. In the approximation $r \equiv \left| \vec{x} \right| \gg \left| \vec{x}_0 \right| \equiv r_0$ we have \begin{equation}
\begin{aligned}
\frac{1}{\left| \vec{x} - \vec{x}_0 \right|} & = \frac{1}{\sqrt{r^2 + r_0^2 - 2 r r_0 \cos \gamma}} \\
& = \frac{1}{r} \LP 1 - \frac{2r_0 \cos \gamma}{r} + \frac{r_0^2}{r^2} \RP^{-\frac{1}{2}} \\
& = \frac{1}{r} \LP 1 - \frac{1}{2} \LP \frac{r_0^2}{r^2} - \frac{2r_0 \cos \gamma}{r} \RP + \frac{3}{8} \LP \frac{r_0^2}{r^2} - \frac{2r_0 \cos \gamma}{r} \RP^2 + \cdots  \RP \\
& = \frac{1}{r} \LP 1 + \frac{r_0}{r} \cos \gamma + \frac{r_0^2}{r^2} \frac{3 \cos^2 \gamma - 1}{2}  + \cdots \RP \\
& = \sum_{n=0}^\infty \frac{r_0^n}{r^{n+1}} P_n(\cos \gamma).
\end{aligned}
\end{equation}
Here we intuitively think of $\vec{x}_0$ as the location of a source and $\vec{x}$ as the location of a distant probe, while $\gamma$ is the (necessarily unsigned) angle between the two position vectors. More generally, given a compactly supported source $\rho: \R^3 \to \R$ for the Poisson equation \begin{equation}
\nabla^2_{\R^3} \phi = \rho
\end{equation}
we can expand in components, organizing our result in according to scaling w.r.t. $r$, to find \begin{equation}
\begin{aligned}
\phi(\vec{x}) & = \int_{\R^3} d^3 \vec{x}' \frac{\rho(\vec{x}')}{|\vec{x} - \vec{x}'|} \\
& = \int_{\R^3} d^3 \vec{x}' \rho(\vec{x}') \frac{1}{r} \LP 1 - \frac{2r' \cos \gamma}{r} + \frac{{r'}^2}{r^2} \RP^{-\frac{1}{2}} \\
& = \int_{\R^3} d^3 \vec{x}' \rho(\vec{x}') \frac{1}{r} \LP 1 + \frac{{r'}^2 - 2 \vec{x} \cdot \vec{x}'}{r^2} \RP^{-\frac{1}{2}} \\
& = \int_{\R^3} d^3 \vec{x}' \rho(\vec{x}') \frac{1}{r} \LP 1 - \frac{1}{2 r^2} \LP {r'}^2 - 2 \vec{x} \cdot \vec{x}' \RP + \frac{3}{8r^4} \LP {r'}^2 - 2 \vec{x} \cdot \vec{x}' \RP^2 + \cdots \RP \\
& = \int_{\R^3} d^3 \vec{x}' \rho(\vec{x}') \frac{1}{r} \LP 1 + \frac{\vec{x} \cdot \vec{x}'}{r^2} - \frac{{r'}^2}{2r^2}  + \frac{3 (\vec{x} \cdot \vec{x}')^2}{2 r^4} + \mc{O} \LP \frac{1}{r^3} \RP \RP \\
& = \int_{\R^3} d^3 \vec{x}' \rho(\vec{x}')  \LP \frac{1}{r} + \frac{x_i' x_i}{r^3}  + \frac{\LP 3 x'_i x'_j  - {r'}^2 \delta_{ij} \RP x_i x_j}{2 r^5} + \mc{O} \LP \frac{1}{r^4} \RP \RP \\
& \equiv \frac{Q}{r} + \frac{Q_i x_i}{r^3} + \frac{Q_{ij} x_i x_j}{2 r^5} + \frac{Q_{ijk} x_i x_j x_k}{3! r^7} + \cdots.
\end{aligned}
\end{equation}
Notice that this implicit definition of that $Q's$ in terms of their results after contraction against $x$ strings leaves exactly enough ambiguity to WLOG take the Q's to be totally symmetric; we will take this as a further part of their definition. As a sanity check, we can explicitly compute the Euclidean Laplacian of this expression. In fact all terms exhibit the same behavior here, as exemplified by the fourth one: \begin{equation}
\begin{aligned}
\nabla^2_{\R^3} \frac{Q_{ijk} x_i x_j x_k}{3! r^7} & = \delta_{ab} \PA{^2}{x_a x_b} \LP \frac{Q_{ijk} x_i x_j x_k}{3! r^7} \RP \\
& =  \delta_{ab} \PA{}{x_a} \LP \frac{Q_{bjk} x_j x_k}{2! r^7}  - \frac{7 Q_{ijk} x_i x_j x_k x_b}{3! r^9} \RP \\
& = \delta_{ab} \LP \frac{Q_{ba k} x_k}{r^7} - \frac{7 Q_{bjk} x_j x_k x_a}{2! r^9} - \frac{7 Q_{ajk} x_j x_k x_b}{2! r^9} - \frac{7 Q_{ijk} x_i x_j x_k \delta_{ab}}{3! r^9} + \frac{63 Q_{ijk} x_i x_j x_k x_a x_b}{3! r^{11}} \RP \\
& = \frac{Q_{aa k} x_k}{r^7} - \frac{7 Q_{ajk} x_j x_k x_a}{2! r^9} - \frac{7 Q_{ajk} x_j x_k x_a}{2! r^9} - \frac{7 Q_{ijk} x_i x_j x_k}{2! r^9} + \frac{63 Q_{ijk} x_i x_j x_k}{3! r^9} \\
& = \frac{Q_{aak} x_k}{r^7}.
\end{aligned}
\end{equation}
In general, we have \begin{equation}
\begin{aligned}
\nabla^2_{\R^3} \phi(\vec{x}) & = \delta^{ab} \PA{^2}{x_a x_b} \LP \frac{Q}{r} + \frac{Q_i x_i}{r^3} + \frac{Q_{ij} x_i x_j}{2 r^5} + \frac{Q_{ijk} x_i x_j x_k}{3! r^7} + \cdots \RP \\
& = \frac{Q_{aa}}{r^5} + \frac{Q_{aai}x_i}{r^7} + \frac{Q_{aaij} x_i x_j}{2! r^9} + \cdots + \frac{Q_{aa i_3 \cdots i_{\ell}} x_{i_3} \cdots x_{i_\ell}}{(\ell - 2)! r^{2 \ell + 1}} + \cdots .
\end{aligned}
\end{equation}
If we are being careful this calculation has a restricted domain of validity. First, we have quietly assumed that $\vec{x} \neq 0$, since more generally \begin{equation}
\nabla^2 \LP \frac{1}{r} \RP = - 4 \pi \delta^{(3)}(\vec{x}).
\end{equation}
More interestingly, we know that $\phi(\vec{x})$ is harmonic away from the compact-by-assumption support of $\rho$, which implies that all of the traces must vanish. \begin{equation}
Q_{aa i_3 \cdots i_\ell} = 0.
\end{equation}
But this implies our series expansion for $\nabla^2_{\R^3} \phi(\vec{x})$ vanishes uniformly. So it must be that the domain of validity of this series expansion (which we generically expect to be limited to sufficiently large $r$ anyway for convergence reasons) is at most the exterior of the support of $\rho$, since inside this support we generally have $\nabla^2_{\R^3} \phi(\vec{x}) = \rho(\vec{x}) \neq 0$. So our series expansion holds exactly, but only for sufficiently large $r$. \\
\\
\
In fact we never needed to talk about $\rho$; it was just a way of cooking up a function $\phi$ which, at large distances, is harmonic and decaying. Any scalar function which is harmonic and polynomial-decaying at large $r$ can be written \begin{equation}
\phi(\vec{x}) = \sum_{\ell = 0}^\infty \frac{Q_{i_1 \cdots i_\ell} x_{i_1} \cdots x_{i_\ell}}{r^{2 \ell + 1}} \mbox{ at sufficiently large } r
\end{equation}
where \begin{equation}
\nabla_{\R^3} \phi(\vec{x}) = 0 \mbox{ at large $r$} \implies Q_{aa i_3 \cdots i_\ell} = 0.
\end{equation}
Pulling out factor of $r$ as $\vec{x} = r \hat{x}$ lets us map each level-$\ell$ multipole is mapped invertibly to a combination of level-$\ell$ spherical harmonic, and we get the whole nice story above. For example, in the normalization \begin{equation}
Y_{\ell m} = e^{i m \phi} P_\ell ^m(\cos \theta)
\end{equation}
and by comparison with the $\ell = 1$ associated Legendre polynomials \begin{equation}
\begin{aligned}
P_1^{-1} (x) & = \frac{1}{2} (1 - x^2)^{\frac{1}{2}}, \\
P_1^0 (x) & = x, \\
P_1^1 (x) & = - (1 - x^2)^{\frac{1}{2}}
\end{aligned}
\end{equation}
we find \begin{equation}
\begin{aligned}
\frac{Q_i x_i}{r^3} & = \frac{ Q_1 \sin \theta \cos \phi + Q_2 \sin \theta \sin \phi+ Q_3 \cos \theta}{r^2} \\
%& = \frac{Q_1 \sin \theta \frac{e^{i \phi} + e^{- i \phi}}{2}  - i Q_2 \sin \theta \frac{e^{i \phi} - e^{- i \phi}}{2} + Q_3 \cos \theta}{r^2} \\
& = \frac{e^{i \phi} \frac{Q_1  - i Q_2 }{2} \sin \theta + e^{- i \phi} \frac{Q_1  + i Q_2 }{2}\sin \theta + Q_3 \cos \theta}{r^2 } \\
& = \frac{Q_3 P_1^0(\cos \theta) - e^{i \phi} \frac{Q_1 - i Q_2}{2} P_1^1(\cos \theta) + e^{- i \phi} (Q_1 + i Q_2) P_1^{-1} (\cos \theta)}{r^2} \\
& = \frac{(Q_1 + i Q_2) Y_{1 -1} (\theta, \phi) + Q_3 Y_{10} (\theta, \phi) - \frac{Q_1 - i Q_2}{2} Y_{11}(\theta, \phi)}{r^2}.
\end{aligned}
\end{equation}





\bibliography{bibfile}
\bibliographystyle{utphys}

\end{document}


